% page 127 du fac-simile (image p-126 du scan)
% bloc 165.1 x 250.6 mm ; marge gauche 36.9 mm
\pgc{15.240mm}{0.000mm}{
\l{}
\l{}
\l{\cel{52}117.}
\l{}
\l{}
\l{\sou{docar}. (\sou{trans.,}\cel{1}ad). Komunikar ad ulu (cienco, arto) per lecioni}
\l{\cel{3}intersequanta. - L.}
\l{}
\l{\sou{docento}. En la universitati di Germania, profesoro da qua la}
\l{\cel{3}lecioni pagesas, ne da la stato, ma da la studenti ipsa.-DEFR.}
\l{}
\l{\sou{docila}.\cel{2}Qua havas bona tendenci pri lasar su instruktesar e}
\l{\cel{3}duktesar. - EFIS.}
\l{}
\l{\sou{dodekaedro}. (\sou{geom.}) Poliedro kun dek-e-du edri. - DEFIRS.}
\l{}
\l{\sou{dodekagina}. (\sou{bot.}) Qua havas dek-e-du stamini. - DEFIRS.}
\l{}
\l{\sou{dogo}. (\sou{zool}.)Hundo grosa e kurta, di qua la muzelo esas kurta, la}
\l{\cel{3}maxilo e la mandibulo forta, la labii pendanta. - DF.}
\l{}
\l{\sou{dogano}.\cel{2}Administrerio qua impostas la vari taxita, kande oli}
\l{\cel{3}eniras od ekiras la teritorio. - FIS.}
\l{}
\l{\sou{dogmato}. Kredendajo o punto, en doktrino religiala o filozofiala.}
\l{\cel{3}- DEFIRS.}
\l{}
\l{\sou{dogmatiko}. La teologio qua traktas la dogmato Kristana. - DEFIRS.}
\l{}
\l{\sou{dogro}. Mikra navo pontoza,quan onu uzas por peskar la haringi}
\l{\cel{3}e la makreli. - DEFIS.}
\l{}
\l{\sou{dojo}. Duko elektita, chefo nominala di la republiki olima di}
\l{\cel{3}Venezia e di Genova (Italia). - DEFIR.}
\l{}
\l{\sou{doko}. Maro-baseno por konstruktar o reparar navi ; baseno cir-}
\l{\cel{3}kondata da kayi e warfi por charjar o descharjar navi.- DFR.}
\l{}
\l{\sou{doktoro.}\cel{2}La persono qua obtenis la grado universitatala maxim}
\l{\cel{3}alta en fakultato (teologio, yurocienco, medicino,cienci,}
\l{\cel{3}literaturo). - DEFIRS.}
\l{}
\l{\sou{doktrino}. Ensemblo de nocioni propozata da ulu kom docenda pri}
\l{\cel{3}temo. - DEFIRS.}
\l{}
\l{\sou{dokumento}. Texto skribura, raporto-texto, akto, qua uzesas por}
\l{\cel{3}klarigar fakti historiala, judiciala, e c. - DEFIRS.}
\l{}
\l{\sou{dolo}. (\sou{Yurocienco}). Trompo (per tacar pri ula fakti konocinda da}
\l{\cel{3}la koncernato ; o per kredigar, da ica, fakti qui ne existas)}
\l{\cel{3}qua detrimentas la interesti di ulu. - DFIS.}
\l{}
\l{\sou{dolca}. I. Qua plezas a la gustado pro saporo agreabla. - II.}
\l{\cel{3}(\sou{metaf}.) Qua donas juo subtila. - FIS.}
\l{}
\l{\sou{dolio}. (\sou{tekn.}) Kavajo qua esas ye la extremajo di la instrumento}
\l{\cel{3}specala qua uzesas en la rivetago por impedar la riveto ek-}
\l{\cel{3}batesar per la martelagi. - E.}
}
